\scanpage{354}
\mspage{93}

B) Cas $\tau=\Delta_\varepsilon\Delta_{-\varepsilon}^{\ell}\Delta_\varepsilon$, $\tau'=\mathbb{\Delta}_\varepsilon^2$ ou $\mathbb{\Delta}_\varepsilon$.

Notons que l'on a, pour $\tau,\tau'\unsure{\in|\vdash}\underline{\mathrm{Ty}}$ \unsure{quelconque|quelconques}, un isom. de catégories
\[
\underline{\mathrm{Hom}}(\tau,\tau')\longrightarrow\underline{\mathrm{Hom}}(\tau'{}^\circ,\tau^\circ),
\]
exprimant simplement que
\[
\tau\longmapsto\tau^\circ
\]
est un foncteur \deleted{\unsure{et}} de $\underline{\mathrm{Ty}}$ dans $\underline{\mathrm{Ty}}$, et $\unsure{=|\simeq}$ un automorphisme de cette catégorie. Donc dire que $\underline{\mathrm{Hom}}(\tau,\tau')$ est connexe revient à dire que $\underline{\mathrm{Hom}}(\tau^\circ,\tau'{}^\circ)$ l'est. Donc on peut supposer \unsure{ci-dessous|ci-dessus} que $\varepsilon=+1$ (et pareil dans A !).

\diagram{\begin{tikzcd}[ampersand replacement=\&,column sep=.65em]
\tau:\&{}\arrow[r]\&{}\&{}\arrow[l]\&{}\arrow[l]\&\cdots\&{}\arrow[l]\&{}\arrow[r]\&{}
\arrow[from=1-3,to=1-7,phantom,"{\underbrace{\hspace{3.1cm}}_{\ell}}"']
\end{tikzcd}}
\diagram{\begin{tikzcd}[ampersand replacement=\&,column sep=.8em]
\deleted{\unsure{I}}\tau':\&{}\arrow[r]\&{}\arrow[r]\&{}\&\text{ou}\&{}\arrow[r]\&{}
\end{tikzcd}}

Mais tous les \deleted{foncteurs} morphismes $\tau\mathrel{\unsure{-|\to}}\tau'$ sont triviaux sur l'intervalle critique $I_2$ (rétrograde), ce qui implique que
\[
\underline{\mathrm{Hom}}(\tau,\tau')\xleftarrow{\sim}\underline{\mathrm{Hom}}(\Delta_\varepsilon^2,\tau'),
\]
\deleted{\unsure{en quoi}} \inserted{via} le foncteur de contraction canonique
\[
I_{\unsure{\tau|I}}\mathrel{\unsure{-|\to}}I_{\Delta_\varepsilon^2}=\mathbb{\Delta}^{2}.
\]

Or dans A on a vu que $\underline{\mathrm{Hom}}(\Delta_\varepsilon^2,\tau')$ est connexe (pour $\tau'$ de la forme $\Delta_\varepsilon$ ou $\Delta_\varepsilon^2$).

On a gagné !

% archive sheet number: 353
